\documentclass{situationist-preprint}

\title{The Stability of Knowledge}
\msid{EPI-14-b}
\msfamily{EPI-14}
\msdiscipline{Epistemology}
\mstype{Consequence}
\mskeywords{value of knowledge; stability; Meno; misleading evidence; dogmatism paradox}

\begin{document}
\maketitle

\begin{abstract}
EPI-14-a argues that knowledge is more valuable than mere true belief because it is more stable under rational updating. We draw three consequences. First, the advantage of knowledge grows with the time over which a belief is to be used, and vanishes in the limiting case of a single immediate action, which is the case of the road to Larissa. Second, stability does not license disregarding later evidence, so the account does not lead to the dogmatism paradox. Third, the advantage is greatest where later evidence is most often produced by the facts, which explains why the value of knowledge is felt most in the empirical domain.
\end{abstract}

\section{Three consequences}

By Theorem~3.2 of EPI-14-a, knowledge is retained for longer in expectation than mere justified true belief, under rational responses to later evidence (Definition~2.1). By Principle~P1 there, the value of a true belief accrues over the time it is held and available to guide action. Together these give Corollary~3.3: knowledge is, other things equal, more valuable.

\begin{proposition}[Horizon]\label{p:horizon}
The excess value of knowledge over mere true belief is non-decreasing in the length of time over which the belief is to guide action and inference, and is zero when that time is zero.
\end{proposition}

\begin{proof}[Argument]
By Principle~P1 of EPI-14-a, the value of each belief is the value accrued over the time it is retained, up to the horizon of use. Where the horizon is zero, both are used once and at once, and accrue the same value. As the horizon lengthens, the expected time for which mere true belief has been abandoned while knowledge is retained can only increase, by Theorem~3.2.
\end{proof}

This vindicates Socrates' concession in the \emph{Meno} \citep[97a--98a]{Meno}: for a single journey to Larissa, begun now, true opinion is as good a guide as knowledge. The difference appears only for the traveller who must keep his way over time, as Socrates says, before his opinions run away.

\begin{proposition}[No dogmatism]\label{p:nodog}
That knowledge is stable does not make it rational for a knower to disregard later evidence against what she knows.
\end{proposition}

\begin{proof}[Argument]
Stability, in Definition~2.1 of EPI-14-a, is defined by retention \emph{under rational responses} to later evidence. It is a property that the belief has given rational updating, not a policy that replaces it. The stability of knowledge is the fact that, since later evidence is produced by the facts (Assumption~2.2 there) and a safe belief is connected to those facts (Lemma~3.1), rational responses to it seldom overturn knowledge.
\end{proof}

The dogmatism paradox, which Harman attributes to Kripke, runs: if I know $p$, any evidence against $p$ is misleading, so I may resolve to ignore it \citep{Harman1973}. Proposition~\ref{p:nodog} blocks the step from the value of stability to that resolution. The knower's belief is valuable because it survives the evidence, not because it is shielded from it; a resolution to ignore evidence would make retention independent of the facts, and so remove the connection that explains why knowledge is retained.

\begin{proposition}[Domain]\label{p:domain}
The stability advantage of knowledge over mere true belief is greatest where the evidence received after belief formation is most often produced by the facts on which the belief's basis depends.
\end{proposition}

This follows from the role of Assumption~2.2 in the proof of Theorem~3.2. In empirical inquiry, where the world continually produces evidence about the matters believed, a lucky true belief is soon exposed; in domains where little later evidence arrives, the advantage is smaller. Williamson's case of the burglar who persists in searching a house because he knows, rather than merely truly believes, that there is a diamond in it, illustrates the first kind of domain \citep[ch.~3]{Williamson2000}.

\section{Conclusion}

The stability of knowledge is a temporal and conditional advantage: it grows with the time over which a belief is used, it presupposes rational responsiveness to evidence rather than replacing it, and it is greatest where the facts keep producing evidence. Each feature is predicted by the account of EPI-14-a and matches the Meno's image of opinions that must be tied down if they are to stay.

\begin{thebibliography}{99}\small
\bibitem[Harman(1973)]{Harman1973} Harman, G. (1973). \emph{Thought}. Princeton University Press.
\bibitem[Plato(Meno)]{Meno} Plato. \emph{Meno}. Cited by Stephanus page.
\bibitem[Williamson(2000)]{Williamson2000} Williamson, T. (2000). \emph{Knowledge and Its Limits}. Oxford University Press.
\end{thebibliography}

\end{document}
